% !TeX root = ../../Book2.tex
\chapter{分数积分与 Littlewood–Paley 理论}
\label{b2:ch:19}
\BookChapterNumber{b2:ch:19}
\begin{chapterabstract}
Riesz 势用齐次核提高可积指数，局部平均与远场 Hölder 的平衡产生 Hardy–Littlewood–Sobolev 不等式。Littlewood–Paley 分解则把多尺度频率信息组织成平方函数，并将几乎正交性推广到 \(L^p\)。本章给出两类估计的完整基础证明，再以非齐次块说明 Sobolev、Besov 与 Triebel–Lizorkin 空间之间的联系。
\end{chapterabstract}

\begin{learninggoals}\begin{itemize}
\item\label{b2:item:19-goal-1}计算 Riesz 势的 Fourier 符号和尺度，解释两个端点的失败。
\item\label{b2:item:19-goal-2}由近远场平衡证明 HLS，并推导非端点 Sobolev 不等式。
\item\label{b2:item:19-goal-3}构造二进频率分割，检查重构、环带支集和 Bernstein 常数。
\item\label{b2:item:19-goal-4}用随机符号和统一核界证明平方函数等价。
\item\label{b2:item:19-goal-5}区分空间积分与尺度求和的次序及低频处理。
\end{itemize}\end{learninggoals}

解 Poisson 方程时，右端与解的二阶导数处在相同的 \(L^p\) 范围；若只看解本身，基本解的长程尾部又改变可积指数。Riesz 势将后一个现象独立出来。另一种问题是判断一个函数在不同频率上的大小是否能重构其原来的范数，这时必须保留尺度之间的平方结构。

\ProseParagraphBreak

本章使用\chapref{b2:ch:17}的极大估计和\chapref{b2:ch:18}的标准核算子理论。随机乘子的平方有界性由有限重叠直接得到，空间核界由 Schwartz 衰减求和得到，因此平方函数证明不会循环使用尚未建立的 Mihlin 定理。Loukas Grafakos 的《Modern Fourier Analysis》\cite{Grafakos2014Modern}与 Elias M. Stein 的《Harmonic Analysis: Real-Variable Methods, Orthogonality, and Oscillatory Integrals》\cite{Stein1993Harmonic}提供这些方法的系统背景；本章的具体版本由下列论证建立。

% !TeX root = ../../Book2.tex
\section{Riesz 势}
\label{b2:sec:1901}

卷积核 \(|x|^{\alpha-d}\) 在原点可积，但在无穷远不可积。它改善可积指数而不是保持原来的指数；这种改善由核的齐次尺度决定。

\subsection{分数积分}
\label{b2:subsec:190101}

\begin{definition}\label{b2:thm:19-potential-definition}
设 \(0<\alpha<d\)。对适当的可测 \(f\)，定义\term[riesz-shi]{Riesz 势}[Riesz potential]
\[
 I_\alpha f(x)\coloneqq c_{d,\alpha}\int_{\mathbb R^d}
 |x-y|^{\alpha-d}f(y)\,\mathrm dy,\qquad
 c_{d,\alpha}=\dfrac{\Gamma((d-\alpha)/2)}
 {2^\alpha\pi^{d/2}\Gamma(\alpha/2)}.
\]
\symidx{\(I_\alpha f\)：Riesz 势}
正函数时允许取无穷值；一般函数先要求绝对积分存在。
\end{definition}
此参数 \(\alpha\) 是增加的光滑阶数，不是核中负幂的大小。为了核对 Fourier 常数，利用热核 \(G_t(x)=(4\pi t)^{-d/2}e^{-|x|^2/(4t)}\) 写成
\[
 c_{d,\alpha}|x|^{\alpha-d}
 =\dfrac1{\Gamma(\alpha/2)}
 \int_0^\infty t^{\alpha/2-1}G_t(x)\,\mathrm dt.
\]
换元 \(s=|x|^2/(4t)\) 直接得到右侧的 Gamma 积分。对测试函数的配对可以用截断热积分再取极限；Fourier 热符号给出
\[
 \widehat{I_\alpha f}(\xi)=(2\pi|\xi|)^{-\alpha}\widehat f(\xi).
\]
低频奇性局部可积恰由 \(\alpha<d\) 保证。对非负紧支集输入，空间积分定义与分布定义相同，不需额外的零平均条件。

\subsection{缩放}
\label{b2:subsec:190102}

若 \(f_\lambda(x)=f(\lambda x)\)、\(\lambda>0\)，换元给出
\[
 I_\alpha f_\lambda(x)=\lambda^{-\alpha}(I_\alpha f)(\lambda x).
\]
假如存在尺度一致的 \(L^p\to L^q\) 估计，非零测试函数代入后必须满足
\[
 \dfrac1q=\dfrac1p-\dfrac\alpha d.
\]
这只是必要条件，不给出有界性。特别是同指数 \(L^p\to L^p\) 的全空间估计与齐次尺度不相容；加上低频正则化的 Bessel 势则是另一种非齐次算子。

\ProseParagraphBreak

对 \(f\geq0\) 紧支集且质量非零，远场 \(I_\alpha f(x)\) 与 \(|x|^{\alpha-d}\) 同阶。对 \(p=1\) 的临界目标 \(q=d/(d-\alpha)\)，这个尾部的 \(q\) 次积分正好是对数发散。尺度关系允许的端点仍可能只有弱估计。

\subsection{核估计}
\label{b2:subsec:190103}

下面采用\term[hedberg-guji]{Hedberg 估计}[Hedberg inequality]的下述版本。

\begin{lemma}[Hedberg 估计]\label{b2:thm:19-hedberg}
设 \(1\leq p<d/\alpha\)，\(f\in L^p(\mathbb R^d)\)。则
\[
 I_\alpha|f|(x)\leq C
 \|f\|_p^{\alpha p/d}\cdot
 (Mf(x))^{1-\alpha p/d}
\]
在 \(Mf(x)<\infty\) 处成立。
\end{lemma}
\begin{proof}
对任意 \(R>0\)，近场分成 \(2^{-k-1}R<|x-y|\leq2^{-k}R\) 的环带。每带不超过 \(C(2^{-k}R)^\alpha Mf(x)\)，求和得 \(CR^\alpha Mf(x)\)。远场若 \(p>1\)，Hölder 与极坐标给出 \(CR^{\alpha-d/p}\|f\|_p\)；条件 \(\alpha p<d\) 正好保证远尾核属于 \(L^{p'}\)。当 \(p=1\)，直接用远场核的上界得到相同形式。

若 \(Mf(x)>0\)、\(\|f\|_p>0\)，令 \(R=(\|f\|_p/Mf(x))^{p/d}\)，两个部分同阶。若 \(Mf(x)=0\)，所有中心球上 \(|f|\) 积分为零，递增球穷尽全空间，故 \(f=0\) 几乎处处。零范数同样直接处理。
\end{proof}
近场使用局部平均，远场使用全域可积性。只估计其中一边，不能得到尺度正确的势估计。这个平衡将在临界与非临界指数之间反复出现。

% !TeX root = ../../Book2.tex
\section{Hardy–Littlewood–Sobolev 不等式}
\label{b2:sec:1902}

势估计把几何核的尺度转为函数的可积指数。它也给出 Sobolev 不等式的一种 Fourier 证明，但此证明在 \(p=1\) 的端点需要另换工具。

\subsection{Hardy--Littlewood--Sobolev 不等式}
\label{b2:subsec:190201}

下面采用\term[hardy-littlewood-sobolev-budengshi]{Hardy--Littlewood--Sobolev 不等式}[Hardy--Littlewood--Sobolev inequality]的下述版本。

\begin{theorem}[Hardy--Littlewood--Sobolev]\label{b2:thm:19-hls}
设 \(0<\alpha<d\)、\(1<p<d/\alpha\)，且 \(q=dp/(d-\alpha p)\)。则
\[
 \|I_\alpha f\|_q\leq C_{d,\alpha,p}\|f\|_p.
\]
当 \(p=1\) 时，有 \(I_\alpha:L^1\to L^{d/(d-\alpha),\infty}\) 的弱型估计。
\end{theorem}
\begin{proof}
由 \(|I_\alpha f|\leq I_\alpha|f|\) 和 Hedberg，取 \(q\) 次方后，极大函数的指数为 \((1-\alpha p/d)q=p\)。应用 \(M\) 的强 \(L^p\) 有界性，得到
\[
 \|I_\alpha f\|_q
 \leq C\|f\|_p^{\alpha p/d}\cdot
 \|Mf\|_p^{1-\alpha p/d}
 \leq C\|f\|_p.
\]
这同时保证空间绝对积分几乎处处有限。\(p=1\) 时，势超过 \(t\) 的集合包含于
\(\{Mf>c(t/\|f\|_1^{\alpha/d})^{d/(d-\alpha)}\}\)；
极大弱型给出测度不超过 \(C(\|f\|_1/t)^{d/(d-\alpha)}\)。
\end{proof}
非零非负紧支集输入的远尾否定强 \(L^1\) 端点。另一方面，在 \(p=d/\alpha\) 处，远场 Hölder 的核指数成为临界，证明中的径向积分发散，不能把 \(q=\infty\) 直接代入结论。

\subsection{Sobolev 不等式的 Fourier 证明}
\label{b2:subsec:190202}

\begin{theorem}[Sobolev 不等式的 Fourier 推导]\label{b2:thm:19-sobolev}
设 \(d\geq2\)、\(1<p<d\)，\(u\in C_c^\infty(\mathbb R^d)\)。则
\[
 \|u\|_{dp/(d-p)}\leq C_{d,p}\|\nabla u\|_p.
\]
\end{theorem}
\begin{proof}
在 Fourier 空间中，\(\sum_jR_j\partial_ju\) 的符号是 \(2\pi|\xi|\widehat u(\xi)\)，故
\[
 u=I_1\left(\sum_jR_j\partial_ju\right)
\]
按温和分布成立。括号内函数属于 \(L^p\)，因为 Riesz 变换在这个指数上有界。HLS 给出该势属于 \(L^{dp/(d-p)}\)，分布等式于是也是几乎处处的函数等式。对每个分量估计，再用 \(|\partial_ju|\leq|\nabla u|\)，得到结论。
\end{proof}
这个证明把一阶导数先组合成径向 Fourier 符号，再以一阶势抵消。它不覆盖 \(p=1\)，因为 Riesz 变换一般不在 \(L^1\) 有界；一阶端点 Sobolev 不等式本身成立，但需要几何或坐标积分的另一种证明。对零边界闭包可以通过光滑紧支集逼近传递，任意区域的非零边界函数还需要延拓或额外零阶项。

% 扩写定位：b2:subsec:190202
分数积分的弱端点与一阶 Sobolev 的强端点应当区别。单独应用 \(I_1:L^1\to L^{d/(d-1),\infty}\) 只给出弱结论；梯度各分量来自同一个函数，这个附加结构允许得到强结论。
\begin{theorem}\label{b2:thm:19-sobolev-l1}
设 \(d\geq2\)。对 \(u\in C_c^1(\mathbb R^d)\)，有
\[
 \|u\|_{d/(d-1)}
 \leq\prod_{j=1}^d\|\partial_ju\|_1^{1/d}
 \leq d^{-1}\sum_{j=1}^d\|\partial_ju\|_1.
\]
因此此式通过 \(W^{1,1}\) 的紧支集光滑逼近推广到 \(W^{1,1}(\mathbb R^d)\)。在 \(d=1\) 时相应的结论为 \(\|u\|_\infty\leq\|u'\|_1\)。
\end{theorem}
\begin{proof}
沿第 \(j\) 个坐标用微积分基本定理，有
\[
 |u(x)|\leq F_j(\widehat x_j),\qquad
 F_j(\widehat x_j)=\int_{\mathbb R}|\partial_ju(x_1,\ldots,t,\ldots,x_d)|\,\mathrm dt.
\]
这里 \(\widehat x_j\) 表示删去第 \(j\) 个坐标，因而 \(F_j\) 与该坐标无关。需要的积分不等式是
\[
 \int_{\mathbb R^d}\prod_{j=1}^d F_j(\widehat x_j)^{1/(d-1)}\,\mathrm dx
 \leq\prod_{j=1}^d\left(\int_{\mathbb R^{d-1}}F_j\right)^{1/(d-1)}.
\]
给出归纳证明：\(d=2\) 是 Fubini。对 \(d\geq3\)，先积分 \(x_d\)，对前 \(d-1\) 个因子用指数 \(d-1\) 的 Hölder；得到
\(F_d(x')^{1/(d-1)}\prod_{j<d}G_j(\widehat x'_j)^{1/(d-1)}\)，其中 \(G_j=\int F_j\,\mathrm dx_d\)。
在 \(x'\in\mathbb R^{d-1}\) 上再用指数 \(d-1\) 与 \((d-1)/(d-2)\) 的 Hölder，把 \(F_d\) 分离出来；剩下的积分由 \(d-1\) 维归纳假设估计。所有因子非负，Tonelli 允许先截断再取极限，故不预设乘积可积。

逐点有 \(|u|^{d/(d-1)}\leq\prod_jF_j^{1/(d-1)}\)，代入上述不等式，再取 \((d-1)/d\) 次方即得第一项；第二项是算术--几何平均。对 \(W^{1,1}\) 的逼近序列，差函数也满足所证明的估计，所以它在目标强空间中 Cauchy；其局部 \(L^1\) 极限与原函数一致，得到推广。
\end{proof}
这一证明利用坐标方向，并未把核的 \(L^1\) 映射性质加强。它也说明若只留下一个任意 \(L^1\) 输入，便失去了为什么能够有强端点的理由。对有界区域上的函数，若没有零边界值，还需加上低阶项或减去平均；常函数就是不能省去这一条件的直接反例。

\ProseParagraphBreak

从上述端点还可独立恢复 \(1<p<d\) 的强 Sobolev 估计。对 \(|u|^\gamma\) 使用端点，令 \(\gamma=p(d-1)/(d-p)>1\)。链式法则给出
\[
 \bigl\||u|^\gamma\bigr\|_{d/(d-1)}
 \leq C\gamma\int |u|^{\gamma-1}|\nabla u|
 \leq C\gamma\|u\|_{dp/(d-p)}^{\gamma-1}\|\nabla u\|_p,
\]
因为 \((\gamma-1)p'=dp/(d-p)\)。左端是目标范数的 \(\gamma\) 次方。非零输入可以消去公共因子；零输入直接满足结论。复值输入用模长的 Lipschitz 链式法则，或先以 \((|u|^2+\eta^2)^{1/2}-\eta\) 正则化后取极限。这个坐标证明与前面的 Fourier 势证明得到相同指数，但只有前者已经处理了 \(p=1\)。

\subsection{势估计}
\label{b2:subsec:190203}

势算子也产生双线性估计。若 \(1<p,r<\infty\) 且 \(1/p+1/r=1+\alpha/d\)，则
\[
 \left|\int (I_\alpha f)g\right|
 \leq C\|f\|_p\cdot\|g\|_r.
\]
确实取 \(q=r'\)，HLS 与 Hölder 的指数正好相容。这种形式在能量泛函中常比输出范数更方便；复函数的配对按本卷的双线性分布约定书写，绝对值估计相同。

\ProseParagraphBreak

在临界指数取 \(f(x)=|x|^{-\alpha}(\log(e/|x|))^{-\beta}\mathbf1_{\{|x|<1/2\}}\)。若 \(p=d/\alpha\) 且 \(1/p<\beta\leq1\)，则 \(f\in L^p\)，但 \(I_\alpha f(0)=\infty\)。两个积分分别归结为 \(\int^\infty s^{-\beta p}\,\mathrm ds\) 和 \(\int^\infty s^{-\beta}\,\mathrm ds\)。单点发散不直接否定本性有界；由于核和输入非负，Fatou 还说明 \(I_\alpha f(x)\to\infty\) 当 \(x\to0\)，因此每个足够高的水平都有正测度的超水平集合，确实否定 \(L^\infty\) 目标。

超过维数的梯度可积性则导致连续性，而不再是某个有限指数的改善。下面给出一个直接的局部积分证明，说明临界 \(p=d\) 为什么仍然需要另外处理。
\begin{proposition}\label{b2:prop:19-morrey}
若 \(p>d\) 且 \(u\in W^{1,p}(\mathbb R^d)\)，则存在连续代表，使
\[
 |u(x)-u(y)|\leq C_{d,p}|x-y|^{1-d/p}\|\nabla u\|_p
 \qquad(x,y\in\mathbb R^d).
\]
\end{proposition}
\begin{proof}
先设 \(u\) 光滑。对含 \(x\) 的凸球 \(B\)、半径 \(r\)，沿从 \(x\) 到 \(z\in B\) 的线段积分，再对 \(z\) 平均。以 \(w=x+t(z-x)\) 换元，球包含关系及 \(t^{-d}\) 的积分给出
\[
 |u(x)-u_B|\leq C_d\int_B
 \frac{|\nabla u(w)|}{|x-w|^{d-1}}\,\mathrm dw.
\]
具体地，一个固定 \(w\) 只有在 \(t\geq|w-x|/(2r)\) 时可能出现；换元后系数为
\(|B|^{-1}|w-x|\int_{|w-x|/(2r)}^1t^{-d-1}\,\mathrm dt
\leq C_d|w-x|^{1-d}\)，这补足了平均线段的计算。

Hölder 给出右端至多
\[
 C\|\nabla u\|_{L^p(B)}
 \left(\int_{|v|<2r}|v|^{-(d-1)p'}\,\mathrm dv\right)^{1/p'}
 \leq C_{d,p}r^{1-d/p}\|\nabla u\|_{L^p(B)}.
\]
径向积分收敛恰需 \((d-1)p'<d\)，即 \(p>d\)。对 \(x,y\)，取同时包含二者且半径 \(2|x-y|\) 的球，以同一个平均 \(u_B\) 比较两点便得结论。

一般 Sobolev 输入取全空间光滑逼近。差函数的上述估计与其 \(L^p\) 平均控制，保证逼近列在每个紧球上一致 Cauchy；一致极限为连续代表，且局部 \(L^p\) 收敛使它与 \(u\) 几乎处相同。最后把点差估计传到极限。
\end{proof}
临界 \(p=d\) 使核的 Hölder 积分成为对数发散，本证明不能给出零阶 Hölder 连续性。即使非临界时，连续代表也不是原来任意选定的点值代表；修改零测集上的值可以人为破坏连续性。把 Sobolev 等价类提升为一个连续函数，是上述逼近和局部一致控制共同得到的结论。

\ProseParagraphBreak

局部版本在一个较大球内控制较小球上的点差，常数同时记录两球的间距。若靠近区域边界，需要延拓性质或边界坐标的规则性；单独知道区域内部的梯度可积，并不足以跨过一个断裂或尖细边界比较任意两点。势理论、区域几何与函数空间的区别在这里仍然不能省略。

\ProseParagraphBreak

临界梯度指数也有一个直接的函数反例，无需诉诸势的表示。设 \(d\geq2\)，在半径 \(1/2\) 的球内取
\[
 u(x)=\log\log(e/|x|),
\]
再乘一个在更小球上为一的光滑截断。它趋于无穷，但处于 \(W^{1,d}(\mathbb R^d)\)。函数本身的 \(d\) 次方可积，因为换元 \(v=\log(e/r)\) 后径向积分含有指数衰减因子；梯度的关键积分则为
\[
 \int_0^{1/2}\frac{\mathrm dr}
 {r(\log(e/r))^d}<\infty.
\]
截断的导数只在离原点的环带产生光滑可积项。弱导数公式可以先在去掉小球的区域分部积分，再令小球半径趋零；边界贡献至多为 \(r^{d-1}\log\log(e/r)\)，当 \(d\geq2\) 时趋零。因此没有在原点偷偷增加一个分布点质量。

\ProseParagraphBreak

这个例子排除了 \(W^{1,d}\subset L^\infty\)，也排除了从该范数直接控制连续代表的结论。将 \(u\) 截为 \(\min(u,N)\) 并作紧支撑光滑逼近，可以得到梯度和函数的 \(L^d\) 范数一致有界、而上确界趋无穷的一列光滑函数。这里链式法则使截断梯度不增加；适当小的卷积保留接近最高平台的值。因此失败不限于一个粗糙等价类，也存在于光滑检验函数上。维数一的临界指数为一，绝对连续函数可由导数积分控制上确界，情形不同，不能把这个高维对数反例直接搬到直线上。

% !TeX root = ../../Book2.tex
\section{二进分解}
\label{b2:sec:1903}

势用一个齐次核同时处理所有频率。二进分解则把频率空间拆成比例相同的环带，每带对应一种尺度；它适合记录平滑阶数，并允许不同尺度以平方而非绝对和的方式合成。

\subsection{频率局部化}
\label{b2:subsec:190301}

取实值、径向、光滑函数 \(\chi\)，满足 \(0\leq\chi\leq1\)，在 \(|\xi|\leq1\) 时为一，在 \(|\xi|\geq2\) 时为零，且沿半径单调。令 \(\varphi(\xi)=\chi(\xi)-\chi(2\xi)\)。它支撑于 \(1/2\leq|\xi|\leq2\)，并满足
\[
 \sum_{j\in\mathbb Z}\varphi(2^{-j}\xi)=1\quad(\xi\ne0).
\]
等式由有限和的望远镜消去得到，不能在零频同样成立。频率定位的空间代价是失去紧支集：\(\check\varphi\) 是快速衰减函数，但一般遍布全空间。

\ProseParagraphBreak

一个位置紧支集而又频率紧支集的非零光滑函数不存在。Fourier 积分在紧频带时延拓为整函数，若它在位置上的开区间或开球外为零，唯一性迫使恒为零。后章的波包因此使用快速衰减的空间局部化，而不要求两种支集同时精确紧致。

\subsection{Littlewood--Paley 投影}
\label{b2:subsec:190302}

定义\term[littlewood-paley-touying]{Littlewood--Paley 投影}[Littlewood--Paley projection]
\[
 \Delta_jf\coloneqq\mathcal F^{-1}
 \bigl(\varphi(2^{-j}\,\cdot)\widehat f\bigr)
 =2^{jd}\check\varphi(2^j\,\cdot)*f.
\]
\symidx{\(\Delta_jf\)：Littlewood--Paley 频率投影}
这里的 \(\Delta_j\) 是频率投影记号，不是 Laplace 算子，也不是第一卷的有限差分 \(\Delta_h\)。下标的整数或向量类型区别其含义。
\begin{proposition}\label{b2:thm:19-lp-resolution}
若 \(1<p<\infty\)，则对 \(f\in L^p(\mathbb R^d)\)，有
\[
 f=\lim_{N\to\infty}\sum_{j=-N}^N\Delta_jf
 \quad\text{在 }L^p\text{ 中}.
\]
\end{proposition}
\begin{proof}
有限和乘子为 \(\chi(2^{-N}\xi)-\chi(2^{N+1}\xi)\)。第一项的空间核是积分一的 Schwartz 逼近恒等，故范数趋于 \(f\)。第二项是极低频平均。先对 \(f\in L^1\cap L^p\)，Young 不等式给出其 \(L^p\) 范数不超过 \(C2^{-(N+1)d(1-1/p)}\|f\|_1\)，趋于零。其卷积核的 \(L^1\) 范数一致有界，再以稠密性推广到一般 \(f\in L^p\)。在 \(p=1\) 时指数为零，不能使用这个结论。
\end{proof}

\subsection{二进环带}
\label{b2:subsec:190303}

每个 \(\Delta_jf\) 的频率支撑在半径约 \(2^j\) 的环带，相邻环带只有有限重叠。若 \(|j-k|>2\)，则 \(\Delta_j\Delta_k=0\)。投影并不是两两正交，也一般不是幂等；名称表示光滑频率定位，不能把它们当成 Hilbert 空间的正交投影。
下面采用\term[bernstein-budengshi]{Bernstein 不等式}[Bernstein inequality]的下述版本。

\begin{proposition}[Bernstein 估计]\label{b2:thm:19-bernstein}
若 \(f\in L^p\)、\(1\leq p\leq q\leq\infty\)，且 Fourier 支撑包含于 \(B(0,C2^j)\)，则对非负整数 \(k\)，
\[
 \|D^kf\|_q\leq C_{k,p,q}2^{j(k+d(1/p-1/q))}\|f\|_p.
\]
\end{proposition}
\begin{proof}
取光滑紧支集乘子 \(\eta\)，在该频带上为一。于是 \(D^kf=D^k(2^{jd}\check\eta(2^j\,\cdot))*f\)。令 \(1+1/q=1/p+1/r\)，核的 \(L^r\) 范数由换元等于 \(C2^{j(k+d(1-1/r))}\)。Young 不等式给出公式。恒等式先按分布成立，右侧卷积是所需范数中的函数，故也给出导数的函数代表。
\end{proof}
只知道 Fourier 支集半径还不能得到导数的反向界，因为低频函数可能落在球中心；反向估计需要环带避开零频。

% 扩写定位：b2:subsec:190303
环带避开原点时，一阶导数确实能反向控制函数。设 Fourier 支撑在
\(\{c2^j\leq|\xi|\leq C2^j\}\)，取光滑 \(\eta\) 在此环带上为一且支集仍避开零频。令
\[
 a_\ell(\zeta)=\frac{\eta(\zeta)\zeta_\ell}{2\pi\mathrm i|\zeta|^2}.
\]
则 \(\sum_\ell a_\ell(2^{-j}\xi)\,(2\pi\mathrm i\xi_\ell)=2^j\) 于该支集，故
\[
 f=2^{-j}\sum_{\ell=1}^d a_\ell(2^{-j}D)\partial_\ell f,\qquad
 \|f\|_p\leq C2^{-j}\sum_\ell\|\partial_\ell f\|_p
 \quad(1\leq p\leq\infty).
\]
每个 \(a_\ell\) 光滑紧支集，其逆 Fourier 核可积，缩放保持核 \(L^1\) 范数，因而只用 Young 即可证明这一界。这里并未提前调用后面的 Mihlin 定理；环带上的紧频率符号与全空间零次齐次符号有不同的核可积性。

\ProseParagraphBreak

把这一步重复 \(k\) 次，或者直接以多指标配分 \(|\xi|^{2k}\)，得到
\[
 \|f\|_p\leq C_k2^{-jk}\sum_{|\alpha|=k}\|\partial^\alpha f\|_p.
\]
它与正向 Bernstein 共同说明：一个纯环带函数的一阶导数，尺度就是 \(2^j\)；频率越高，相同函数范数所需的梯度越大。一个低频常数则不能满足这样的反向不等式，因为它的梯度为零。非齐次 Sobolev 范数中的低频 \(L^p\) 项正是为了补回这一信息。

\subsection{近似正交性}
\label{b2:subsec:190304}

在 \(L^2\) 中，由 Plancherel，
\[
 \sum_j\|\Delta_jf\|_2^2
 =\int\left(\sum_j|\varphi(2^{-j}\xi)|^2\right)|\widehat f(\xi)|^2\,\mathrm d\xi.
\]
分割非负且和为一，重叠次数不超过固定整数 \(L\)，故括号在非零频处介于 \(1/L\) 与一之间。这给出平方和与 \(L^2\) 范数的等价，常数依赖所选分割而不是输入。

\ProseParagraphBreak

若一列函数分别具有相互分离的频率支集，则其平方范数可以相加；若支集只有限重叠，就在重叠次数内损失一个常数。这是\term[ji-hu-zhengjiao]{几乎正交性}[almost orthogonality]的最基本版本。对 \(L^p\) 不能直接使用 Plancherel，下一节用随机符号与奇异积分把同样的尺度信息传递到全部开区间指数。

% !TeX root = ../../Book2.tex
\section{平方函数}
\label{b2:sec:1904}

平方函数把每个频带的振幅合成一个位置函数。它不要求频带输出逐点正交；随机符号使抵消的平均能够准确记录平方大小。

\subsection{Littlewood--Paley 平方函数}
\label{b2:subsec:190401}

定义\term[pingfang-hanshu]{平方函数}[square function]
\[
 Sf(x)\coloneqq\left(\sum_{j\in\mathbb Z}|\Delta_jf(x)|^2\right)^{1/2}.
\]
\symidx{\(Sf\)：Littlewood--Paley 平方函数}
级数先以单调极限解释，暂时允许无穷值。取独立符号 \(\varepsilon_j=\pm1\)，两值等概率。对任意有限复数族 \(a_j\)，有 Khintchine 估计
\[
 c_p\left(\sum_j|a_j|^2\right)^{p/2}
 \leq\mathbb E\left|\sum_j\varepsilon_ja_j\right|^p
 \leq C_p\left(\sum_j|a_j|^2\right)^{p/2}.
\]
为看清常数与项数无关，先取实数族。乘积公式和 \(\cosh s\leq e^{s^2/2}\) 给出
\(\mathbb E e^{t\sum\varepsilon_ja_j}\leq e^{t^2\sum a_j^2/2}\)；
优化 \(t\) 得到 Gaussian 尾界，再按层蛋糕积分得到所有有限 \(p\) 的上界。平方期望为 \(A^2=\sum a_j^2\)，四次期望不超过 \(3A^4\)。在 \(\{|X|^2>A^2/2\}\) 上用 Cauchy–Schwarz，得到该集合概率至少 \(1/12\)，遂给出下界。复数族分别处理实部、虚部，其中一个平方和至少总和的一半，常数只多一个固定因子。

\subsection{\texorpdfstring{\(L^p\)}{Lp} 等价性}
\label{b2:subsec:190402}

下面采用\term[littlewood-paley-budengshi]{Littlewood--Paley 不等式}[Littlewood--Paley inequality]的下述版本。

\begin{theorem}[Littlewood--Paley]\label{b2:thm:19-square-function}
对 \(1<p<\infty\)，存在只依赖 \(p,d,\chi\) 的正常数，使
\[
 c_p\|f\|_p\leq\|Sf\|_p\leq C_p\|f\|_p
 \quad(f\in L^p(\mathbb R^d)).
\]
\end{theorem}
\begin{proof}
对有限指标集 \(J\) 和任意符号，令 \(T_\varepsilon=\sum_{j\in J}\varepsilon_j\Delta_j\)。其乘子一致有界，故 \(L^2\) 范数一致有界。核为 \(\sum_{j\in J}\varepsilon_j2^{jd}\check\varphi(2^jx)\)。利用 Schwartz 衰减，按 \(2^j|x|\leq1\) 与 \(>1\) 分开几何级数，分别得到
\[
 |K_\varepsilon(x)|\leq C|x|^{-d},\quad
 |\nabla K_\varepsilon(x)|\leq C|x|^{-d-1},
\]
常数与 \(J,\varepsilon\) 无关。每项有限，支集外表示直接成立，Calderón–Zygmund 定理给出一致 \(L^p\) 界。Khintchine 与 Tonelli 得到有限平方和的上界，再以单调收敛推广到全部指标。

取实偶光滑环带函数 \(\widetilde\varphi\)，在 \(\varphi\) 的支集上为一，记对应投影为 \(\widetilde\Delta_j\)。同一个随机核论证给出 \(\widetilde Sg\) 的上界。对 Schwartz \(f,g\)，频率分割和 Plancherel 给出
\[
 \int f\overline g
 =\sum_j\int\Delta_jf\,\overline{\widetilde\Delta_jg}.
\]
先取有限和再取极限，或者用 \(Sf\cdot\widetilde Sg\) 的可积性控制级数。逐点 Cauchy–Schwarz 和 Hölder 给出配对不超过 \(C_{p'}\|Sf\|_p\cdot\|g\|_{p'}\)。对 \(g\) 取对偶上确界，得到下界；密度与已有上界推广到 \(L^p\)。
\end{proof}
因此有界性不依赖尺度数量，逐带得到的常数不会在无穷求和时积累。

% 扩写定位：b2:subsec:190402
平方函数还给出级数重构的无条件性质。若 \(f\in L^p\)、\(1<p<\infty\)，则
\[
 \sum_{j\in\mathbb Z}c_j\Delta_jf,\qquad |c_j|\leq1,
\]
按任意指标排列均在 \(L^p\) 中收敛，并且其范数至多 \(C_p\|f\|_p\)。这里系数固定后，不能把原级数的极限仍说成 \(f\)；只有全部 \(c_j=1\) 才重构原函数。

\ProseParagraphBreak

证明的要点是控制任何有限尾集 \(J\)，而不是只控制对称截断。对 \(g\in L^{p'}\)，在配对中插入等于一的环带投影，得到
\[
 \left|\int\sum_{j\in J}c_j\Delta_jf\,\overline g\right|
 \leq \left\|\left(\sum_{j\in J}|\Delta_jf|^2\right)^{1/2}\right\|_p
       \|\widetilde Sg\|_{p'}.
\]
因而有限尾和的范数由对应平方尾控制。\(Sf\in L^p\)，且平方尾在逐点上趋零，支配收敛说明这种控制一致趋零。于是所有排列都是 Cauchy；完整平方函数上界给出统一范数。对于满足同样有限重叠环带支撑的任意函数族 \(u_j\)，同一配对也给出
\[
 \left\|\sum_j u_j\right\|_p
 \leq C_p\left\|\left(\sum_j|u_j|^2\right)^{1/2}\right\|_p
\]
的有限和版本，只要所插入的环带投影在各个支撑上为一。这正是合成算子的有界性；单独知道每个 \(u_j\) 的 \(L^p\) 大小并不能保证右端可积。

\ProseParagraphBreak

具体地，若各个频块在同一个位置区域上叠加，平方求和比绝对求和节省了大量信息；若各块分别位于远离的区域，则空间积分又会记录它们的位置分离。Besov 范数先取每个频块的全空间大小，已经不再保留第二种位置关系；Triebel--Lizorkin 范数先逐点合成，仍然能区别这两种情形。这解释了为什么两种范数在 \(p=q\) 时相同，而在其他指数中并没有简单的逐项互换。

\subsection{向量值估计}
\label{b2:subsec:190403}

下面记录后续空间估计使用的向量版本。若 \(A_j\) 的核 \(k_j\) 都满足统一大小和一阶差分条件，且每个 \(A_j\) 的 \(L^2\) 范数一致有界，则
\[
 \left\|\left(\sum_{j=1}^N|A_jf_j|^2\right)^{1/2}\right\|_p
 \leq C_p\left\|\left(\sum_{j=1}^N|f_j|^2\right)^{1/2}\right\|_p,
 \quad1<p<\infty,
\]
常数与 \(N\) 无关。这里核条件指对角矩阵核的算子范数，即 \(\sup_j|k_j|\) 和对应差分。该版本的证明可直接重做上一章分解：对向量 \(F=(f_j)\) 的 Euclidean 模作二进选块，块内减去向量平均。好函数模受平均模控制，坏函数的向量积分为零；尾估计只用矩阵核范数，平方估计则按分量相加。弱型插值与 Euclidean 对偶给出同一 \(L^p\) 结论，没有维数为 \(N\) 的常数。

\ProseParagraphBreak

对 \(A_j=\Delta_j\)，核的缩放使这些统一界成立。这项向量估计比“每个投影分别有界”更强：后者只能逐项相加，可能产生随频带数增长的常数。它还说明把随机化用于一个向量族时，应检查其共同核界，而不只检查每个标量算子。

\subsection{Sobolev 范数}
\label{b2:subsec:190404}

在平方可积的 Sobolev 空间中，频率权重直接给出阶数。对 \(s\in\mathbb R\)，以
\[
 \|f\|_{H^s}^2\coloneqq
 \int(1+4\pi^2|\xi|^2)^s|\widehat f(\xi)|^2\,\mathrm d\xi
\]
定义非齐次 \(H^s\)。令 \(P_0=\chi(D)\)，\(P_j=\Delta_j\) 对 \(j\geq1\)。则
\[
 \|f\|_{H^s}^2\asymp
 \|P_0f\|_2^2+\sum_{j\geq1}2^{2js}\|P_jf\|_2^2.
\]
这里 \(\asymp\) 表示双向不等式，常数只依赖固定参数和分割。证明在每个环带上比较 \(1+4\pi^2|\xi|^2\) 与 \(2^{2j}\)，再用有限重叠的平方上下界；低频由 \(P_0\) 单独处理。整数阶 \(s=k\geq0\) 时，这与弱导数定义的 \(H^k\) 范数等价，但系数和多重指标计数可能不同，不能写成数值完全相等。

\ProseParagraphBreak

对一般 \(p\)，空间平方合成应放在 \(L^p\) 范数内部。把 \(\sum2^{2js}\|P_jf\|_p^2\) 当成同一个范数会改变求和与积分的次序。这两种次序正是下一节两族函数空间的区别。

\ProseParagraphBreak

一个简单的频率分割还给出平方 Sobolev 范数的插值不等式。设 \(s>d/2\)、\(u\in\mathcal S\)。按半径 \(R\) 分开 Fourier 积分，Cauchy--Schwarz 给出
\[
 \|\widehat u\|_1\leq C_dR^{d/2}\|u\|_2
 +C_{d,s}R^{d/2-s}\|(-\Delta)^{s/2}u\|_2.
\]
第一项使用球体积，第二项使用 \(\int_{|\xi|>R}|\xi|^{-2s}\mathrm d\xi<\infty\)，严格条件 \(s>d/2\) 就用于此处。反演控制 \(\|u\|_\infty\)，取 \(R^s\) 与两个平方范数的比值相平衡，得到
\[
 \|u\|_\infty\leq C
 \|u\|_2^{1-d/(2s)}
 \|(-\Delta)^{s/2}u\|_2^{d/(2s)}.
\]
若最高阶范数为零，Fourier 支集只能落在原点，平方输入因而为零，可以单独处理。

\ProseParagraphBreak

再把 \(L^2\) 与 \(L^\infty\) 用 Hölder 插值，可得对 \(2\leq q\leq\infty\)
\[
 \|u\|_q\leq C
 \|u\|_2^{1-\theta}\|(-\Delta)^{s/2}u\|_2^\theta,
 \qquad \theta=\frac d s\left(\frac12-\frac1q\right).
\]
这是一种明确范围的 Gagliardo--Nirenberg 不等式。此处采用的最高阶大于半维数，所以证明可以直接通过绝对 Fourier 反演完成；更低阶或更广指数版本需要前面 Sobolev 嵌入及其他插值，不能把该推导的条件自行删去。

\ProseParagraphBreak

低频平方大小与高频导数大小承担不同职责。若只保留高频项，就会失去函数整体振幅和大尺度分布的控制；若只保留低频项，又不能约束尖峰。选取平衡半径正是让这两种信息在同一尺度相遇，这与分数积分证明中的近场、远场平衡具有相同的优化结构。

% !TeX root = ../../Book2.tex
\OptionalSection{Besov 与 Triebel–Lizorkin 空间}{b2:sec:1905}

齐次频带只记录非零频率，因而会丢掉多项式和低频信息。本节使用非齐次分解，在一个低频块之外只保留正整数尺度，以明确空间定义和边界问题的范围。

\subsection{二进定义}
\label{b2:subsec:190501}

令 \(P_0=\chi(D)\)、\(P_j=\Delta_j\) 对 \(j\geq1\)，则 \(\sum_{j\geq0}P_j=I\)。对 \(s\in\mathbb R\)、\(1\leq p<\infty\)、\(1\leq q\leq\infty\)，定义
\[
 \|f\|_{B^s_{p,q}}
 \coloneqq\left\|\{2^{js}\|P_jf\|_p\}_{j\geq0}\right\|_{\ell^q},
 \qquad
 \|f\|_{F^s_{p,q}}
 \coloneqq\left\|\left(\sum_{j\geq0}|2^{js}P_jf|^q\right)^{1/q}\right\|_p,
\]
当 \(q=\infty\) 时用上确界。这些条件在温和分布中定义相应空间。每个频带输出是光滑函数，但其全空间 \(L^p\) 范数可能无穷；只有满足上述有限性条件的分布才属于该空间。

\subsection{Sobolev 空间 的重新解释}
\label{b2:subsec:190502}

平方函数定理给出 \(F^0_{p,2}=L^p\) 且范数等价，\(1<p<\infty\)；上一节给出 \(B^s_{2,2}=F^s_{2,2}=H^s\)。前者要将低频块加入齐次平方函数，使用低频卷积的 \(L^p\) 有界性及分割重构；后者直接由 Plancherel。

\ProseParagraphBreak

一般 \(p\) 的 Bessel 势空间定义为 \(J^sf\in L^p\)，其中 \(J^s\) 的符号是 \((1+4\pi^2|\xi|^2)^{s/2}\)。把它与 \(F^s_{p,2}\) 识别，需要统一控制环带上的归一化符号及其逆。这个乘子步骤将在\chapref{b2:ch:22}建立；此处的平方可积识别已经独立完成，不把未证明的乘子结论藏进定义。

\subsection{Besov 空间}
\label{b2:subsec:190503}

满足第一项范数有限的空间称为\term[besov-kongjian]{Besov 空间}[Besov space]。参数 \(q\) 衡量不同尺度的幅值如何求和。对 \(p=q\)，Tonelli 使 \(B^s_{p,p}\) 与 \(F^s_{p,p}\) 的定义范数相同；当两指数不同时，求和顺序一般不能交换。

\ProseParagraphBreak

例如取频率分别位于半径 \(2^j\) 的不交小球的函数 \(f_j\)，归一化为 \(\|f_j\|_2=1\)。函数 \(\sum_ja_j2^{-js}f_j\) 的 \(B^s_{2,q}\) 大小与 \(\{a_j\}\) 的 \(\ell^q\) 大小可比，有限重叠只改变固定常数。因此 \(q=1,2,\infty\) 分别要求绝对和、平方和或统一有界，不能单由光滑指数 \(s\) 判断空间关系。

% 扩写定位：b2:subsec:190503
对 \(0<s<1\)，Besov 的差分描述可以在有限指数中完整验证。以下使用非齐次范数，因而低频项必须保留。
\begin{proposition}\label{b2:prop:19-besov-difference}
设 \(0<s<1\)、\(1\leq p<\infty\)。对 \(f\in L^p(\mathbb R^d)\)，有
\[
 \|f\|_{B^s_{p,p}}\asymp
 \|f\|_p+
 \left(\int_{\mathbb R^d}
 \frac{\|f(\,\cdot+h)-f\|_p^p}{|h|^{d+sp}}\,\mathrm dh\right)^{1/p}.
\]
因此此处的 Besov 空间与按 Gagliardo 差分定义的分数 Sobolev 空间相同，范数等价。
\end{proposition}
\begin{proof}
先取光滑输入，记 \(u_j=P_jf\)、\(a_j=2^{js}\|u_j\|_p\)。Bernstein 与平移不变性给出
\[
 \|u_j(\,\cdot+h)-u_j\|_p
 \leq C\min\{1,2^j|h|\}\|u_j\|_p.
\]
对 \(2^{-k-1}<|h|\leq2^{-k}\)、\(k\geq0\)，于是
\[
 2^{ks}\|f(\,\cdot+h)-f\|_p
 \leq C\sum_{j\leq k}2^{-(k-j)(1-s)}a_j
       +C\sum_{j>k}2^{-(j-k)s}a_j.
\]
两边对尺度取 \(\ell^p\)；右侧是与两个可和几何序列卷积，离散 Young 不等式给出 \(\ell^p\) 界。每个环带的 \(|h|^{-d}\mathrm dh\) 积分是固定常数，故这就是小步长差分积分的上界。大步长部分由 \(\|f(\cdot+h)-f\|_p\leq2\|f\|_p\) 与 \(s>0\) 控制。

反向利用 \(j\geq1\) 时 \(\int\check\varphi=0\)。置 \(k_j(h)=2^{jd}\check\varphi(2^jh)\)，则
\[
 P_jf(x)=\int k_j(h)[f(x-h)-f(x)]\,\mathrm dh.
\]
Hölder 的加权形式给出
\(\|P_jf\|_p^p\leq\|k_j\|_1^{p-1}\int|k_j(h)|\|f(\cdot-h)-f\|_p^p\,\mathrm dh\)，在 \(p=1\) 时同式直接成立。Schwartz 衰减和几何级数说明
\[
 \sum_{j\geq1}2^{jsp}|k_j(h)|\leq C|h|^{-d-sp},\qquad h\ne0.
\]
将 \(2^j|h|\) 在一的两侧分开求和即可；高频侧取衰减次数大于 \(d+sp\)。Tonelli 得到高频 Besov 范数的上界，低频 \(P_0\) 由 Young 控制。

还需控制右式中的 \(\|f\|_p\)。有限 Besov 范数保证
\(\sum_{j\geq1}\|P_jf\|_p=\sum2^{-js}a_j\leq C_s\|\{a_j\}\|_{\ell^p}\)。
因此重构级数绝对收敛于 \(L^p\)，并按温和分布分割恒等式等于 \(f\)。前面的估计可由频率截断和 Fatou 推广到全部输入，而不依赖未证明的点态级数收敛。
\end{proof}
范围 \(0<s<1\) 的两端都有原因：\(s>0\) 使高频重构及大步长积分收敛；\(s<1\) 使低于步长尺度的几何序列可和。若 \(s\geq1\)，应使用更高阶差分，不能原样沿用一阶差分公式。在 \(s=1\) 时，一阶差分积分的临界径向因子是 \(\mathrm dr/r\)，一般非零梯度已使它发散。

\ProseParagraphBreak

例如，对一个有光滑边界的有界区域 \(E\)，平移使指示函数的变化发生在厚度约 \(|h|\) 的边界层，因此
\(\|\mathbf1_E(\cdot+h)-\mathbf1_E\|_p^p\leq C|h|\)；
在边界具有一块非平坦或平坦图形且取一组横穿该图形的方向时，也有相应下界 \(c|h|\)。径向差分积分于是归结为 \(\int_0^1r^{-sp}\,\mathrm dr\)，给出 \(\mathbf1_E\in B^s_{p,p}\) 当且仅当 \(sp<1\)。这与跳跃穿过一条余维一边界的事实相符：无论区域多么光滑，函数本身的跳跃仍然限制其分数光滑性。

% 扩写定位：b2:subsec:190503
Besov 的绝对尺度求和在非线性问题中有一个直接用途。Bernstein 给出
\(\|P_jf\|_\infty\leq C2^{jd/p}\|P_jf\|_p\)；
因此 \(B^{d/p}_{p,1}\) 连续嵌入有界连续函数空间。每个块连续，且其上确界范数绝对可和，故重构级数一致收敛。这里第二指数 \(1\) 正好允许绝对求和；只知道尺度序列有界并不足以宣布一致收敛。
\begin{proposition}\label{b2:prop:19-besov-algebra}
设 \(1\leq p<\infty\)、\(s>d/p\)。则 \(B^s_{p,1}\) 是一个乘法代数，并有
\[
 \|fg\|_{B^s_{p,1}}\leq C\|f\|_{B^s_{p,1}}\|g\|_{B^s_{p,1}}.
\]
\end{proposition}
\begin{proof}
置 \(f_k=P_kf\)、\(g_l=P_lg\)。两者在 \(L^p\) 和 \(L^\infty\) 中均绝对求和，所以双级数表示的乘积收敛于 \(L^p\)，且与逐点乘积相同。乘积 \(f_kg_l\) 的频率支撑包含在半径 \(C2^{\max\{k,l\}}\) 的球中；这是 Fourier 卷积的支集包含关系，先对光滑输入验证，再在所用范数下逼近。因而 \(P_j(f_kg_l)=0\) 当 \(j>\max\{k,l\}+C_0\)。

对 \(k\geq l\)，Young 和 Hölder 给出
\[
 \sum_j2^{js}\|P_j(f_kg_l)\|_p
 \leq C_s2^{ks}\|f_k\|_p\|g_l\|_\infty.
\]
这里 \(\sum_{0\leq j\leq k+C_0}2^{js}\leq C_s2^{ks}\)，使用了 \(s>0\)。对 \(l>k\) 交换两个函数。把所有指标相加得到
\[
 \|fg\|_{B^s_{p,1}}\leq C_s
 \left(\sum_k2^{ks}\|f_k\|_p\right)\left(\sum_l\|g_l\|_\infty\right)
 +C_s\left(\sum_k\|f_k\|_\infty\right)
 \left(\sum_l2^{ls}\|g_l\|_p\right).
\]
由先前嵌入，两个上确界范数和都受各自 \(B^s_{p,1}\) 范数控制。频率截断与 Fatou 把有限和估计推广到全部输入。
\end{proof}
这段证明已经出现低频与高频相互作用的基本分类。若 \(k\) 明显大于 \(l\)，乘积的尺度仍约为 \(2^k\)，低频因子可以放在上确界范数；若 \(k,l\) 接近，高频可能相互抵消并产生低频，但求和 \(j\leq k+C_0\) 已把这些输出一并控制。只保留高低项而丢掉相近频带，会遗漏非线性乘积的重要部分。

\ProseParagraphBreak

例如，在这个空间中 \(f^2,f^3\) 都保持相同光滑阶数，其范数由 \(\|f\|_{B^s_{p,1}}^2,\|f\|_{B^s_{p,1}}^3\) 控制。它并不说明任意低光滑输入的乘积都在同一空间，更不说明任意两个分布可以相乘。可相乘的关键是已经建立的有界代表与绝对重构；在较低光滑阶数上，应重新分配两因子的可积指数或使用更精细的频率乘积分解。

\subsection{Triebel--Lizorkin 空间}
\label{b2:subsec:190504}

满足第二项范数有限的空间称为\term[triebel-lizorkin-kongjian]{Triebel--Lizorkin 空间}[Triebel--Lizorkin space]。它先在每个位置合成频率块，再作空间积分。平方函数的 \(F^s_{p,2}\) 属于这一族，而把每块 \(L^p\) 范数先算出来的 Besov 形式属于另一族。

\ProseParagraphBreak

在所给有限指数范围中，换用另一套光滑分割得到等价范数：每个新块与有限个相邻旧块相交；旧分割恒等式把新块写成这些项的和，归一化卷积核一致 \(L^1\) 有界。Besov 情形直接在序列上用有限重叠；Triebel–Lizorkin 情形还需相应的向量卷积估计。平方序列 \(q=2\) 已由上一节矩阵核论证覆盖；一般 \(q\) 的完整分割独立性属于该理论的进一步部分，不能仅靠标量 Young 不等式宣布。Hans Triebel 的《Theory of Function Spaces》\cite{Triebel1983Theory}系统处理这些更广指数。

% 扩写定位：b2:subsec:190504
求和与积分的顺序也产生一组基本嵌入。对有限 \(p,q\)，置 \(a_j(x)=2^{js}|P_jf(x)|\)。若 \(p\geq q\)，Minkowski 在指数 \(p/q\) 中给出
\[
 \left\|\sum_j a_j^q\right\|_{p/q}
 \leq\sum_j\|a_j^q\|_{p/q},
\]
即 \(\|f\|_{F^s_{p,q}}\leq\|f\|_{B^s_{p,q}}\)；而逐点 \(\ell^p\) 范数不超过 \(\ell^q\) 范数，给出
\(\|f\|_{B^s_{p,p}}\leq\|f\|_{F^s_{p,q}}\)。
若 \(p\leq q\)，对向量值积分应用 Minkowski，得到
\[
 \left(\sum_j\left[\int a_j^p\right]^{q/p}\right)^{p/q}
 \leq\int\left(\sum_j a_j^q\right)^{p/q};
\]
另一方向由逐点 \(\ell^q\leq\ell^p\) 得到。因此两种情形可合写为
\[
 B^s_{p,\min\{p,q\}}\ \subset\ F^s_{p,q}\
 \subset\ B^s_{p,\max\{p,q\}}.
\]
当 \(q=\infty\) 时，把有限序列上确界取极限也得到同样结论。这里显示的是连续嵌入；是否严格还取决于空间、参数和函数族，不能从一个方向的范数不等式直接宣布所有情况都严格。

\ProseParagraphBreak

改变可积指数时，Bernstein 还给出
\[
 B^s_{p,q}\subset B^{s-d(1/p-1/r)}_{r,q},
 \qquad1\leq p\leq r\leq\infty.
\]
证明只需在每个块上使用 \(\|P_jf\|_r\leq C2^{jd(1/p-1/r)}\|P_jf\|_p\)，再按所给 \(q\) 求和。低频块也满足同样界，因为其频率尺度为一。这个嵌入损失的光滑阶数完全由空间缩放决定；若不损失该阶数，缩放一个固定高频函数便会使候选不等式的两侧具有不同增长幂。

\subsection{迹理论 的联系}
\label{b2:subsec:190505}

频率块也解释边界迹的阶数损失。若一个 \(d\) 维块的频率大小为 \(2^j\)，在法向上把厚度 \(2^{-j}\) 的区域压到边界，\(L^p\) 大小会损失 \(2^{j/p}\)。因此一阶空间的迹阶数为 \(1-1/p\)，一般 \(s\) 阶的自然目标为 \(s-1/p\)。

\ProseParagraphBreak

这只是尺度说明，不能代替迹算子的构造、分割独立性或边界坐标变换估计。通常的非齐次迹结论需要 \(s>1/p\) 和适当规则边界；临界等号处还要记录序列指数，不能直接给出所有空间统一的边界值。第一卷的弱导数迹理论与这里的频率语言描述相同的阶数损失，但数学对象及端点必须逐项对应。本附节不以一个尺度计算冒称完整的 Besov 迹定理。


\begin{chaptersummary}
势的尺度关系规定可积指数，极大函数把这个必要关系变成非端点的充分估计。二进分割的重构需要分别处理高低频，平方合成的常数则来自统一核界和有限重叠。低频块、序列指数和空间积分的顺序决定不同光滑空间的边界；把这些区别保留下来，才可以准确使用后续乘子和迹理论。
\end{chaptersummary}

\BookExercises

% 自拟：围绕本章已经建立的工具，改变条件或组合方法。
\begin{exercise}\label{b2:ex:19-new-1}
设存在非零测试函数上的尺度一致估计 \(\|I_\alpha f\|_q\leq C\|f\|_p\)。证明 \(1/q=1/p-\alpha/d\)，说明为何它不能独自证明端点有界。
\end{exercise}
\begin{solution}
代入 \(f(\lambda x)\) 后，两侧因子分别为 \(\lambda^{-\alpha-d/q}\) 和 \(\lambda^{-d/p}\)。若幂不同，令 \(\lambda\) 趋零或无穷就与固定常数矛盾。非负紧支集输入在 \(p=1,q=d/(d-\alpha)\) 时仍有不可积的临界尾部，故必要尺度不保证强估计。
\end{solution}

% 自拟：围绕本章已经建立的工具，改变条件或组合方法。
\begin{exercise}\label{b2:ex:19-new-2}
令 \(f=\mathbf1_{B(0,1)}\)。给出 \(|x|\geq2\) 时 \(I_\alpha f(x)\) 的双向幂界，证明它不属于临界强目标空间。
\end{exercise}
\begin{solution}
对球内 \(y\)，有 \(|x|/2\leq|x-y|\leq3|x|/2\)。负幂比较给出 \(c|x|^{\alpha-d}\leq I_\alpha f(x)\leq C|x|^{\alpha-d}\)。取 \(q=d/(d-\alpha)\) 后径向积分为常数倍的 \(\int_2^\infty r^{-1}\,\mathrm dr\)，发散。球内核因 \(\alpha>0\) 可积，失败来自远尾而不是原点。
\end{solution}

% 自拟：围绕本章已经建立的工具，改变条件或组合方法。
\begin{exercise}\label{b2:ex:19-new-3}
对临界 \(p=d/\alpha\)，验证正文的对数输入属于 \(L^p\) 而势不是本性有界，并补明由单点发散到正测度超水平集的步骤。
\end{exercise}
\begin{solution}
换元 \(s=\log(e/r)\) 后，输入积分为 \(\int^\infty s^{-\beta p}\,\mathrm ds\)，收敛；中心势为 \(\int^\infty s^{-\beta}\,\mathrm ds\)，发散。对任意 \(x_n\to0\)，除 \(y=0\) 外核逐点趋于中心核；非负 Fatou 给出 \(\liminf_nI_\alpha f(x_n)=\infty\)。若某高度没有一个邻域上的下界，就能取趋零而势不超过该高度的点列，矛盾。因此超水平集合包含小球。
\end{solution}

% 自拟：围绕本章已经建立的工具，改变条件或组合方法。
\begin{exercise}\label{b2:ex:19-new-4}
设 \(0<\alpha,\beta\)、\(\alpha+\beta<d\)，对非负紧支集光滑 \(f\) 证明 \(I_\alpha I_\beta f=I_{\alpha+\beta}f\)，解释参数限制。
\end{exercise}
\begin{solution}
热积分中 \(e^{t\Delta}e^{s\Delta}=e^{(t+s)\Delta}\)。Tonelli 换序，令 \(r=t+s\)、\(u=t/(t+s)\)，Jacobi 行列式为 \(r\)。内层为 \(\int_0^1u^{\alpha/2-1}(1-u)^{\beta/2-1}\,\mathrm du=\Gamma(\alpha/2)\Gamma(\beta/2)/\Gamma((\alpha+\beta)/2)\)，可由两个 Gamma 积分的同一换元证明。余下正是 \(\alpha+\beta\) 阶热表示。该和小于 \(d\) 保证空间核远场的卷积积分及低频齐次分布的所需范围。
\end{solution}

% 自拟：围绕本章已经建立的工具，改变条件或组合方法。
\begin{exercise}\label{b2:ex:19-new-5}
证明频率支集在环带 \(c2^j\leq|\xi|\leq C2^j\) 时，\(\|f\|_2\leq C2^{-j}\|\nabla f\|_2\)。说明球支集不能代替环带。
\end{exercise}
\begin{solution}
Plancherel 给出 \(\|\nabla f\|_2^2=4\pi^2\int|\xi|^2|\widehat f|^2\geq4\pi^2c^22^{2j}\|f\|_2^2\)。若只在球内，可把频率支集继续集中到任意小球，归一化 \(\|f\|_2=1\) 后梯度范数趋零，反向常数不可能统一。
\end{solution}

% 自拟：围绕本章已经建立的工具，改变条件或组合方法。
\begin{exercise}\label{b2:ex:19-new-6}
取 \(f\in\mathcal S\) 且 \(\int f\ne0\)。证明齐次部分和 \(\sum_{j=-N}^N\Delta_jf\) 不可能在 \(L^1\) 中趋于 \(f\)。
\end{exercise}
\begin{solution}
每个 \(\Delta_jf\) 的积分为零，因为乘子在零频为零，所以所有有限和积分为零。若 \(L^1\) 收敛到 \(f\)，积分作为连续线性泛函也收敛，迫使 \(\int f=0\)，矛盾。这不否定 \(1<p<\infty\) 的重构，积分不是这些全空间 \(L^p\) 上的连续泛函。
\end{solution}

% 自拟：围绕本章已经建立的工具，改变条件或组合方法。
\begin{exercise}\label{b2:ex:19-new-7}
设每点至多 \(L\) 个非负频率分割函数非零，且其和为一。证明对应 \(L^2\) 平方函数常数可取 \(L^{-1/2}\) 与一；解释为何通常不是等距。
\end{exercise}
\begin{solution}
对非零频率，Cauchy–Schwarz 给出 \(1\leq L\sum_j\varphi_j^2\)，非负性又给出 \(\sum_j\varphi_j^2\leq(\sum_j\varphi_j)^2=1\)。乘 \(|\widehat f|^2\) 积分。若两个分割在正测度频带同时严格介于零和一，平方和严格小于一；取变换支撑在该处的非零输入，就不是等距。
\end{solution}

% 自拟：围绕本章已经建立的工具，改变条件或组合方法。
\begin{optionalexercise}\label{b2:ex:19-new-8}
取正 \(s\) 和互不相交、位于半径 \(2^j\) 的光滑频率小球，在其中选取 \(\|f_j\|_2=1\)。用 \(f=\sum_{j\geq1}j^{-a}2^{-js}f_j\) 比较 \(B^s_{2,q}\) 与 \(H^s\) 的序列条件。
\end{optionalexercise}
\begin{solution}
正 \(s\) 使原函数的平方范数和收敛，故 \(f\in L^2\)。选小球使其仅与固定数目的相邻环带相交，并至少有一个分割因子一致下界；有限重叠与 Plancherel 给出 Besov 范数等价于 \(\{j^{-a}\}\) 的 \(\ell^q\) 范数，有限 \(q\) 时条件为 \(aq>1\)，无穷指数时为 \(a\geq0\)。\(H^s\) 条件为 \(2a>1\)。选择不同 \(a,q\) 可以分别满足或违反两种条件，展示第二指数的实质作用。
\end{solution}
